\documentclass[../main.tex]{subfiles}
\zexternaldocument*{category_theory}
\begin{document}
\chapter{Homological algebra}

\section{(More) category theory}
\subsection{Additive categories}
\begin{definition}[Additive \nameref{def:category}] \label{def:additive-category}
  A \nameref{def:category} \(\mathsf{A}\) is \textbf{additive} if it verifies:
  \begin{enumerate}[label=A.\arabic*), ref=A.\arabic*]
    \item\label{axiom:additive-zero} There exists a \nameref{def:zero-object} \(0 \in \mathsf{A}\).
    \item\label{axiom:additive-products} \(\mathsf{A}\) has finite \namerefs{def:categorical-product}.
    \item\label{axiom:additive-hom} For every pair of objects \(X, Y \in \mathsf{A}\) the set \(\Hom_{\mathsf{A}}(X, Y)\) is endowed with an abelian group structure and the composition maps are bilinear with respect to the group operation.
  \end{enumerate}
\end{definition}
\SetPluralName{def:additive-category}{Additive \namerefs{def:category}}

\begin{lemma} \label{lemma:products-coproducts-additive}
  Finite \namerefs{def:categorical-product} and \namerefs{def:categorical-coproduct} coincide in an \nameref{def:additive-category}.
\end{lemma}
\begin{proof}
  Consider \(X_{1}, X_{2} \in \mathsf{A}\). We only need to check that \(X_{1} \times X_{2} \cong X_{1} \amalg X_{2}\) since the case for a finite collection of elements follows directly by induction.
  To do this, we will see that \(X_{1} \times X_{2}\) satisfies the universal property of the \nameref{def:categorical-coproduct}.

  Consider the following diagram
  \begin{center}
    \begin{tikzcd}
      X_{1} \arrow[dd, "\id"'] \arrow[dr, "{(\id, 0)}"] &       & X_{2} \arrow[dd, "\id"] \arrow[dl, "{(0, \id)}"'] \\
        & \arrow[dl, "p_{1}"] X_{1} \times X_{2} \arrow[dr, "p_{2}"'] &   \\
      X_{1} &       & X_{2}
    \end{tikzcd}
  \end{center}

  Note that for any pair of morphisms \((\alpha_{1}, \beta_{1}), (\alpha_{2}, \beta_{2}): R \to X_{1} \times X_{2}\) we have that
  \begin{equation*}
    (\alpha_{1}, \beta_{1}) + (\alpha_{2}, \beta_{2}) = (\alpha_{1} + \alpha_{2}, \beta_{1}, \beta_{2})
  \end{equation*}
  Therefore if \(p_{1}, p_{2}\) are two projections on \(X_{1} \times X_{2}\) whose sum is the identity we have that
  \begin{equation*}
    (\id, 0) \circ p_{1} + (0, \id) \circ p_{2} = \id_{X_{1} \times X_{2}}
  \end{equation*}

  Consider now two morphisms \(f_{i}: X_{i} \to Q\).
  Then its sum \(\phi = f_{1} \circ p_{1} + f_{2} \circ p_{2}\) produces the following diagram
  \begin{center}
    \begin{tikzcd}
      X_{1} \arrow[dd, "\id"'] \arrow[dr, "{(\id, 0)}"] &       & X_{2} \arrow[dd, "\id"] \arrow[dl, "{(0, \id)}"'] \\
        &  X_{1} \times X_{2} \arrow[dd, "\phi"] &   \\
        X_{1} \arrow[dr, "f_{1}"'] &       & \arrow[dl, "f_{2}"] X_{2} \\
              &   Q   &
    \end{tikzcd}
  \end{center}
  Furthermore, this is unique.
  Indeed, consider \(\phi': X_{1} \times X_{2} \to Q\) another morphism filling the diagram. Then
  \begin{align*}
    \phi &= \phi \circ \id_{X_{1} \times X_{2}} \\
      &= \phi \circ ( (\id, 0) \circ p_{1} + (0, \id) \circ p_{2} ) \\
      &= \phi \circ (id, 0) \circ p_{1} + \phi \circ (0, \id) \circ p_{2} \\
      &= f_{1} \circ p_{1} + f_{2} \circ p_{2} \\
      &= \phi' \circ (\id, 0) + \phi' \circ (0, \id) \\
      &= \phi' \circ ( (\id, 0) \circ p_{1} + (0, \id) \circ p_{2} ) \\
      &= \phi' \circ \id_{X_{1} \times X_{2}} \\
      &= \phi'
  \end{align*}

  Therefore \(X_{1} \times X_{2}\) verifies the universal property of \(X_{1} \amalg X_{2}\).
\end{proof}

\begin{notation}
  If \(\{X_{i}\}_{i=1}^{n}\) is a finite set of objects in an \nameref{def:additive-category} we will denote its \nameref{def:categorical-product} (and therefore by \autoref{lemma:products-coproducts-additive} its \nameref{def:categorical-coproduct}) by \(\bigoplus_{i=1}^{n} X_{i}\).
\end{notation}

\begin{lemma} \label{lemma:mono-epi-additive}
  A morphism \(\varphi: A \to B\) in an \nameref{def:additive-category} is a \hyperref[def:mono-epi-morphism]{monomorphism} if and only if for every morphism \(\xi: Z \to A\)
  \begin{equation*}
    \varphi \circ \xi = 0 \implies \xi = 0
  \end{equation*}
  Dually \(\varphi\) is an \hyperref[def:mono-epi-morphism]{epimorphism} if and only if for every morphism \(\xi: B \to Z\)
  \begin{equation*}
    \xi \circ \varphi = 0 \implies \xi = 0
  \end{equation*}
\end{lemma}

\begin{definition}[Additive \nameref{def:functor}] \label{def:additive-functor}
  Let \(\mathsf{A}, \mathsf{B}\) be two \namerefs{def:additive-category}. A \nameref{def:functor} \(F: \mathsf{A} \to \mathsf{B}\) is \textbf{additive} if it preserves the abelian group structure of all Hom-sets.
\end{definition}

\begin{definition}[Kernel] \label{def:categorical-kernel}
  Let \(f: A \to B\) be a morphism in an \nameref{def:additive-category} \(\mathsf{A}\).
  A morphism \(\iota: K \to A\) is a \textbf{kernel} of \(f\) if \(f \circ \iota = 0\) and for every morphism \(g: Z \to A\) such that \(f \circ g = 0\) there exists a unique \(\tilde{g}:Z \to K\) making the following diagram commute
  \begin{center}
    \begin{tikzcd}
      Z \arrow[rr, "g"'] \arrow[dr, dotted, "\tilde{g}"'] \arrow[rrr, bend left, "0"] &   & A \arrow[r, "f"'] & B\\
        & K \arrow[ur, "\iota"'] &   &
    \end{tikzcd}
  \end{center}
\end{definition}

\begin{definition}[Cokernel] \label{def:categorical-cokernel}
  Let \(f: A \to B\) be a morphism in an \nameref{def:additive-category} \(\mathsf{A}\).
  A morphism \(\pi: B \to C\) is a \textbf{cokernel} of \(f\) if \(\pi \circ f = 0\) and for every morphism \(g: B \to Z\) such that \(\pi \circ g = 0\) there exists a unique \(\tilde{g}:C \to Z\) making the following diagram commute
  \begin{center}
    \begin{tikzcd}
        &   & C \arrow[dr, dotted, "\tilde{g}"] &   \\
      A \arrow[r, "f"] \arrow[rrr, bend right, "0"'] & B \arrow[rr, "g"] \arrow[ur, "\pi"] &   & Z
    \end{tikzcd}
  \end{center}
\end{definition}

\begin{lemma}
  In an \nameref{def:additive-category} \namerefs{def:categorical-kernel} are \hyperref[def:mono-epi-morphism]{monomorphisms} and \namerefs{def:categorical-cokernel} are \hyperref[def:mono-epi-morphism]{epimorphisms}.
\end{lemma}
\begin{proof}
  We will only show that \namerefs{def:categorical-cokernel} are \hyperref[def:mono-epi-morphism]{epimorphisms}. The statement about \namerefs{def:categorical-kernel} follows analogously.

  Let \(\varphi: A \to B\) be a morphism and let \(\coker \varphi: B \to C\) be its \nameref{def:categorical-cokernel}.
  Consider another morphism \(\gamma: C \to Z\) such that \(\gamma \circ \coker \varphi = 0\).
  Since \((\gamma \circ \coker \varphi) \circ \varphi = 0\) by definition of the \nameref{def:categorical-cokernel} we have
  \begin{center}
    \begin{tikzcd}
      A \arrow[r, "\varphi"] & B \arrow[dr, "\coker \varphi"'] \arrow[rr, "\gamma \circ \coker \varphi = 0"] &   & Z \\
        &   & C \arrow[ur, bend left, "\gamma"'] \arrow[ur, dotted, bend right, "\exists!"'] &
    \end{tikzcd}
  \end{center}
  Therefore since \(\gamma \circ \coker \varphi = 0 = 0 \circ \coker \varphi\) uniqueness forces \(\gamma = 0\) and thus by \autoref{lemma:mono-epi-additive} \(\varphi\) is an \hyperref[def:mono-epi-morphism]{epimorphism}.
\end{proof}

\begin{lemma} \label{lemma:partial-sequence-additive}
  Let \(\varphi: A \to B\) be a morphism in an \nameref{def:additive-category}. Then \(\varphi\) is a \hyperref[def:mono-epi-morphism]{monomorphism} if and only if \(0 \to A\) is its \nameref{def:categorical-kernel} and \(\varphi\) is an \hyperref[def:mono-epi-morphism]{epimorphism} if and only if \(B \to 0\) is its \nameref{def:categorical-cokernel}.
\end{lemma}
\begin{proof}
  We will do the \nameref{def:categorical-kernel} one.
  \begin{itemize}
    \item[\(\implies\)] Suppose \(\varphi\) is a \hyperref[def:mono-epi-morphism]{monomorphism}.

          If \(\xi: Z \to A\) is a morphism such that \(\varphi \circ \xi = 0\) then by \autoref{lemma:mono-epi-additive} \(\xi = 0\).
          In particular \(\xi\) factors uniquely through \(0 \to A\), i.e: \(0 \to A\) is the \nameref{def:categorical-kernel} of \(\varphi\).
    \item[\(\impliedby\)] Suppose now that \(0 \to A\) is the \nameref{def:categorical-kernel} of \(\varphi\).

          Consider a morphism \(\xi: Z \to A\) such that \(\varphi \circ \xi = 0\).
          By the \nameref{def:categorical-kernel} properties we have the following diagram
          \begin{center}
            \begin{tikzcd}
              Z \arrow[rr, "\xi"'] \arrow[rrr, bend left, "0"] \arrow[dr, dotted, "\exists!"'] & & A \arrow[r, "\varphi"'] & B \\
              & 0 \arrow[ur]
            \end{tikzcd}
          \end{center}
          Therefore by uniqueness \(\xi = 0\), making \(\varphi\) a \hyperref[def:mono-epi-morphism]{monomorphism} by \autoref{lemma:mono-epi-additive}.
  \end{itemize}
\end{proof}


\subsection{Abelian categories}
\begin{definition}[Abelian \nameref{def:category}] \label{def:abelian-category}
  An \nameref{def:additive-category} \(\mathsf{A}\) is \textbf{abelian} if it verifies
  \begin{enumerate}
    \item[A.4)]\label{axiom:abelian-kernel} \namerefs{def:categorical-kernel} and \namerefs{def:categorical-cokernel} exist in \(\mathsf{A}\).

          Furthermore every \hyperref[def:mono-epi-morphism]{monomorphism} is the \nameref{def:categorical-kernel} of some morphism and every \hyperref[def:mono-epi-morphism]{epimorphism} is the \nameref{def:categorical-cokernel} of some morphism.
  \end{enumerate}
\end{definition}

\begin{lemma} \label{lemma:kernel-of-cokernel-abelian}
  In an \nameref{def:abelian-category} \(\mathsf{A}\) every \nameref{def:categorical-kernel} is the \nameref{def:categorical-kernel} of its \nameref{def:categorical-cokernel} and every \nameref{def:categorical-cokernel} is the \nameref{def:categorical-cokernel} of its \nameref{def:categorical-kernel}.
\end{lemma}
\begin{proof}
  Lets check \namerefs{def:categorical-cokernel} this time.

  Let \(\coker \varphi: A \to B\) be the \nameref{def:categorical-cokernel} of a morphism \(\varphi: Z \to A\).
  Since \(\mathsf{A}\) is \hyperref[def:abelian-category]{abelian} \(\varphi\) has a kernel \(\ker \varphi: K \to A\).
  Note that \(Z \to A \to B\) is 0, therefore by the \nameref{def:categorical-kernel} universal property we have
  \begin{center}
    \begin{tikzcd}
      Z \arrow[rr, "\varphi"] \arrow[dr, dotted] & & A \arrow[r, "\coker \varphi"'] & B \\
      & K \arrow[ur, "\ker \varphi"']
    \end{tikzcd}
  \end{center}

  Consider a morphism \(A \to C\) such that \((K \to A \to C) = 0\).
  Now using the \nameref{def:categorical-cokernel} universal property we obtain
  \begin{center}
    \begin{tikzcd}
      &&& C & \\
      Z \arrow[rr, "\varphi"] \arrow[dr, dotted] \arrow[urrr, bend left, dotted, "0"] & & A \arrow[rr, "\coker \varphi"'] \arrow[ur] & & \arrow[ul, dotted] B \\
      & K \arrow[ur, "\ker \varphi"']
    \end{tikzcd}
  \end{center}
  Therefore \(\coker \varphi\) satisfies the universal property of the \nameref{def:categorical-cokernel} of \(\ker \varphi: K \to A\).
\end{proof}

\begin{lemma}
  Let \(\varphi: A \to B\) be a morphism in an \nameref{def:abelian-category} \(\mathsf{A}\).
  If \(\varphi\) is both a \hyperref[def:mono-epi-morphism]{monomorphism and epimorphism} then \(\varphi\) is an \nameref{def:categorical-isomorphism}.
\end{lemma}
\begin{proof}
  By \autoref{lemma:partial-sequence-additive} the \nameref{def:categorical-kernel} of \(\varphi\) is \(0 \to A\) and the \nameref{def:categorical-cokernel} is \(B \to 0\).
  Furthermore by \autoref{lemma:kernel-of-cokernel-abelian} \(\varphi\) is the \nameref{def:categorical-cokernel} of \(0 \to A\) and the \nameref{def:categorical-kernel} of \(B \to 0\).

  Consider the following diagram
  \begin{center}
    \begin{tikzcd}
        &   & B \arrow[d, "\id"] &   \\
      0 \arrow[r] & A \arrow[r, tail, two heads, "\varphi"] & B \arrow[r] & 0
    \end{tikzcd}
  \end{center}
  Since \((B \to B \to 0) = 0\) and \(\varphi\) is the \nameref{def:categorical-kernel} of \(B \to 0\) we have
  \begin{center}
    \begin{tikzcd}
        &   & B \arrow[d, "\id"] \arrow[dl, dotted, "\psi"'] &   \\
      0 \arrow[r] & A \arrow[r, tail, two heads, "\varphi"] & B \arrow[r] & 0
    \end{tikzcd}
  \end{center}
  Therefore since \(\varphi \psi = \id\) we have shown that \(\varphi\) has a right-inverse.

  Consider now the diagram
  \begin{center}
    \begin{tikzcd}
                  &   A & &   \\
      0 \arrow[r] & A \arrow[u, "\id"] \arrow[r, tail, two heads, "\varphi"] & B \arrow[r] & 0
    \end{tikzcd}
  \end{center}
  Since \((0 \to A \to A) = 0\) and \(\varphi\) is the \nameref{def:categorical-cokernel} of \(0 \to A\) we have
  \begin{center}
    \begin{tikzcd}
                  &   A & &   \\
      0 \arrow[r] & A \arrow[u, "\id"] \arrow[r, tail, two heads, "\varphi"] & B \arrow[ul, dotted, "\eta"'] \arrow[r] & 0
    \end{tikzcd}
  \end{center}
  Therefore since \(\eta \varphi = \id\) we have shown that \(\varphi\) has a left-inverse.
\end{proof}



\begin{theorem}[Freyd-Mitchell theorem] \label{thm:freyd-mitchell}
  Let \(\mathsf{A}\) be a small \nameref{def:abelian-category}. Then there exists a \hyperref[def:functor-types]{fully faithfull} \nameref{def:exact-functor} \(\mathsf{A} \to R\text{-}\mathsf{Mod}\) for a suitable (possibly non-commutative) ring \(R\), where \(R\text{-}\mathsf{Mod}\) is the category of left \(R\)-modules.
\end{theorem}


\section{Complexes and (co)homology}

From now on, we will be working in a fixed \nameref{def:abelian-category} \(\mathsf{A}\).
By the \nameref{thm:freyd-mitchell} we can work with objects and morphisms using set-theoretic tools when needed.

\begin{definition}[Complex] \label{def:complex}
  A sequence of objects and morphisms in an \nameref{def:additive-category} \(\mathsf{A}\)
  \begin{center}
    \begin{tikzcd}
      \cdots \arrow[r, "d_{i-2}"] & C_{i-1} \arrow[r, "d_{i-1}"] & C_{i} \arrow[r, "d_{i}"] & C_{i+1} \arrow[r, "d_{i+1}"] & \cdots
    \end{tikzcd}
  \end{center}
  is a \textbf{chain complex} if for every \(i \quad d_{i} \circ d_{i-1} = 0\).
  We denote a chain complex as \((C_{\bullet}, d_{\bullet})\) or simply as \(C_{\bullet}\).

  By reversing the order, we say that a sequence
  \begin{center}
    \begin{tikzcd}
      \cdots \arrow[r, "d^{i+2}"] & C^{i+1} \arrow[r, "d^{i+1}"] & C^{i} \arrow[r, "d^{i}"] & C^{i-1} \arrow[r, "d^{i-1}"] & \cdots
    \end{tikzcd}
  \end{center}
  is a \textbf{cochain complex} if for every \(i \quad d_{i} \circ d_{i+1} = 0\).
  We denote a cochain complex as \((C^{\bullet}, d^{\bullet})\) or simply as \(C^{\bullet}\).
\end{definition}

\begin{remark}
  In purelly algebraic terms, chains are cochains are \textit{equivalent}.
  We will adopt the cochain notation.
\end{remark}


\begin{definition}[Exact complex] \label{def:exact-complex}
  A \nameref{def:complex} in an \nameref{def:abelian-category}
  \begin{center}
    \begin{tikzcd}
      \cdots \arrow[r, "d^{i+2}"] & C^{i+1} \arrow[r, "d^{i+1}"] & C^{i} \arrow[r, "d^{i}"] & C^{i-1} \arrow[r, "d^{i-1}"] & \cdots
    \end{tikzcd}
  \end{center}
  is \textbf{exact at \(C^{i}\)} (or simply at position \(i\)) if
  \begin{equation*}
    \coker d^{i+1} \circ \ker d^{i} = 0
  \end{equation*}

  A complex is \textbf{exact} if it is exact at every position.
\end{definition}

\begin{definition}
  Let \(C^{\bullet}\) be a \nameref{def:complex} in an \nameref{def:abelian-category} \(\mathsf{A}\).
  We define the \textbf{\(i\)-th cohomology group} of \(C^{\bullet}\) as
  \begin{equation*}
    \operatorname{H}^{i}(C^{\bullet}) = \quot{\ker d^{i}}{\im d^{i-1}}
  \end{equation*}
\end{definition}


\end{document}
